% !TEX program = xelatex
% !TEX encoding = UTF-8 Unicode
%
% Research-oriented resume for Zhijun Yang (Chinese version).
% Layout adapted from Jake's Resume:
% https://github.com/jakegut/resume

\documentclass[a4paper,11pt]{article}

\usepackage{fontspec}
\setmainfont[
  BoldFont=texgyretermes-bold.otf,
  ItalicFont=texgyretermes-italic.otf,
  BoldItalicFont=texgyretermes-bolditalic.otf,
  Ligatures=NoCommon
]{texgyretermes-regular.otf}
\usepackage{xeCJK}
\setCJKmainfont[
  AutoFakeBold=2,
  AutoFakeSlant=.2
]{Droid Sans Fallback}
\setCJKsansfont[
  AutoFakeBold=2,
  AutoFakeSlant=.2
]{Droid Sans Fallback}
\XeTeXlinebreaklocale "zh"
\XeTeXlinebreakskip=0pt plus 1pt
\usepackage[
  a4paper,
  left=0.70in,
  right=0.70in,
  top=0.58in,
  bottom=0.58in
]{geometry}
\usepackage{enumitem}
\usepackage{tabularx}
\usepackage{titlesec}
\usepackage[hidelinks]{hyperref}
\usepackage{xcolor}

\hypersetup{
  pdftitle={杨志军——中文求职简历},
  pdfauthor={杨志军},
  pdfsubject={计算机系统与分布式系统方向求职简历},
  pdfkeywords={分布式系统, 分离式内存, CXL, RDMA, 资源池化}
}

\pagestyle{empty}
\setlength{\parindent}{0pt}
\setlength{\parskip}{0pt}
\setlength{\tabcolsep}{0pt}
\urlstyle{same}
\raggedright

\titleformat{\section}
  {\large\bfseries}
  {}
  {0pt}
  {}
  [\vspace{0.08em}\titlerule]
\titlespacing*{\section}{0pt}{0.78em}{0.42em}

\newcommand{\educationentry}[4]{%
  \begin{tabularx}{\textwidth}{@{}Xr@{}}
    \textbf{#1} & #2 \\
    \textit{#3} & \textit{#4}
  \end{tabularx}%
}

\newcommand{\datedentry}[2]{%
  \begin{tabular*}{\textwidth}{@{\extracolsep{\fill}}lr}
    #1 & #2
  \end{tabular*}%
}

\newcommand{\resumelink}[2]{\href{#1}{[#2]}}

\begin{document}

\begin{center}
  {\Huge\bfseries 杨志军}\\[6pt]
  {\normalsize
    华中科技大学计算机系统结构博士研究生
  }\\[3pt]
  {\small
    \href{tel:+8613599296060}{135-9929-6060}
    \enspace$\vert$\enspace
    \href{mailto:cszjyang@hust.edu.cn}{cszjyang@hust.edu.cn}
    \enspace$\vert$\enspace
    \href{https://cszjyang.github.io/}{cszjyang.github.io}
    \enspace$\vert$\enspace
    \href{https://github.com/cszjyang}{github.com/cszjyang}
  }
\end{center}
\vspace{-0.55em}

\section{教育经历}

\educationentry
  {华中科技大学}
  {中国，武汉}
  {计算机系统结构，博士研究生}
  {2023 年 9 月 -- 2028 年 6 月（预计）}

\vspace{0.15em}
\datedentry
  {\textit{导师：}\href{https://csyhua.github.io/csyhua/index.html}{华宇教授}}
  {}

\vspace{0.30em}
\educationentry
  {华中科技大学}
  {中国，武汉}
  {计算机科学与技术，工学学士}
  {2019 年 9 月 -- 2023 年 6 月}

\section{研究方向}

分布式系统、分离式内存、基于 CXL/RDMA 的系统，以及资源解耦与池化。

\section{论文发表}

\begin{enumerate}[
  leftmargin=1.55em,
  label={\textup{[\arabic*]}},
  itemsep=0.42em,
  topsep=0.05em,
  parsep=0pt,
  partopsep=0pt
]

\item
  \textbf{Zhijun Yang}, Yu Hua, Ming Zhang, Menglei Chen, Xumin Chen, and Aoyang Tong.
  ``XTRA: Unifying Cache Coherence and Concurrency Control for Distributed Transactions in a CXL Pod.''
  \textit{32nd ACM Symposium on Operating Systems Principles (SOSP '26),} 即将发表。
  \resumelink{https://github.com/cszjyang/xtra}{代码}

\item
  \textbf{Zhijun Yang}, Yu Hua, Ming Zhang, Menglei Chen, and Yixiao Wang.
  ``FORGE: Mitigating Synchronization Amplification for Memory-Disaggregated Caching Systems.''
  \textit{20th USENIX Symposium on Operating Systems Design and Implementation (OSDI '26),}
  pp.~977--995, 2026.
  \textbf{获 Jay Lepreau 最佳论文奖提名。}
  \resumelink{https://www.usenix.org/conference/osdi26/presentation/yang-zhijun}{论文}
  \resumelink{https://github.com/cszjyang/forge}{代码}

\item
  Menglei Chen, Yu Hua, \textbf{Zhijun Yang}, Yixiao Wang, and Xumin Chen.
  ``CloverRec: Cost-Efficient Disaggregated Processing-in-Memory Systems for Deep Recommendation Inference.''
  \textit{25th USENIX Conference on File and Storage Technologies (FAST '27),} 即将发表。
  \resumelink{https://github.com/LighT-chenml/CloverRec}{代码}

\item
  Ming Zhang, Yu Hua, and \textbf{Zhijun Yang}.
  ``Motor: Enabling Multi-Versioning for Distributed Transactions on Disaggregated Memory.''
  \textit{18th USENIX Symposium on Operating Systems Design and Implementation (OSDI '24),}
  pp.~801--819, 2024.
  \resumelink{https://www.usenix.org/conference/osdi24/presentation/zhang-ming}{论文}
  \resumelink{https://github.com/minghust/motor}{代码}

\end{enumerate}

\section{学术服务}

\datedentry
  {EuroSys 2026 影子程序委员会委员（春季及秋季）}
  {2026}
\vspace{0.24em}
\datedentry
  {USENIX ATC 2025 工件评审委员会委员}
  {2025}

\section{荣誉与奖励}

\datedentry
  {SOSP 2026 学生奖学金（Student Scholarship）}
  {2026}
\vspace{0.20em}
\datedentry
  {OSDI 2026 学生资助（USENIX Student Grant）}
  {2026}
\vspace{0.20em}
\datedentry
  {华中科技大学优秀研究生}
  {2024}
\vspace{0.20em}
\datedentry
  {华中科技大学科技创新奖学金}
  {2024}
\vspace{0.20em}
\datedentry
  {华中科技大学优秀本科毕业生}
  {2023}
\vspace{0.20em}
\datedentry
  {新加坡国立大学计算机学院暑期项目优秀表现}
  {2021}

\end{document}
